% Generated by scripts/assemble.py; corrections belong in source/corrections.json.
\examhead{44}{上海交通大学 2026 年数学分析试卷}{46}
% p047-b14; source PDF page 47
\par\Needspace{6\baselineskip}
\blockanchor{p047-b14}
\question{1}计算下列极限\par
\nopagebreak[4]
% p047-b15; source PDF page 47
\blockanchor{p047-b15}
\displayformula{\lim _ {x \rightarrow 0} \frac {1 + \cos x - \cos x \cos (2 x) - \sqrt [ 3 ]{\cos x}}{\ln (1 - x) + \ln (1 + x)}.}
% p047-b16; source PDF page 47
\par\Needspace{6\baselineskip}
\blockanchor{p047-b16}
\question{2}设 $u(x, y)$ 在 $\mathbb{R}^2$ 上具有二阶连续偏导数，并且满足\par
\nopagebreak[4]
% p047-b17; source PDF page 47
\blockanchor{p047-b17}
\displayformula{\frac {\partial^ {2} u}{\partial x ^ {2}} - \frac {\partial^ {2} u}{\partial y ^ {2}} = 0.}
% p047-b18; source PDF page 47
\blockanchor{p047-b18}
若 $u(x,2x) = x,u_x'(x,2x) = x^2$ ，求 $u_{xx}^{\prime \prime}(x,2x),u_{xy}^{\prime \prime}(x,2x),u_{yy}^{\prime \prime}(x,2x).$\par
% p047-b19; source PDF page 47
\par\Needspace{6\baselineskip}
\blockanchor{p047-b19}
\question{3}设 $f(x, y, z)$ 是 $\Omega = \{(x, y, z) \in \mathbb{R}^3 \mid x^2 + y^2 + z^2 \leq 1\}$ 上的黎曼可积函数, $p \in \mathbb{R}$, 若 $|f(x, y, z)|$ 在原点附近有正下界, 试求广义积分\par
\nopagebreak[4]
% p047-b20; source PDF page 47
\blockanchor{p047-b20}
\displayformula{I = \iiint_ {\Omega} \frac {f (x , y , z) \ln \sqrt {x ^ {2} + y ^ {2} + z ^ {2}}}{(x ^ {2} + y ^ {2} + z ^ {2}) ^ {p}} \mathrm{d} x \mathrm{d} y \mathrm{d} z}
% p047-b21; source PDF page 47
\blockanchor{p047-b21}
收敛时 p 的取值范围.\par
% p047-b22; source PDF page 47
\blockanchor{p047-b22}
\question{4}设  $f(x)$  在区间  $[0, a]$  上有二阶连续导数， $f'(0) = 1, f''(0) \neq 0$ ，且  $0 < f(x) < x, x \in (0, a)$ ，任取  $x_{1} \in (0, a)$ ，令  $x_{n+1} = f(x_{n}), n = 1, 2, \cdots$ ，证明：数列  $\{x_{n}\}$  和  $\{nx_{n}\}$  都收敛.\par
% p048-b00; source PDF page 48
\blockanchor{p048-b00}
\question{5}证明: 若函数 $f(x)$ 在 $\mathbb{R}$ 上有定义, 并且满足: $\forall x \in \mathbb{R}, \exists \delta_x > 0$, 当 $x_1, x_2 \in (x - \delta_x, x + \delta_x)$ 且 $x_1 < x < x_2$ 时, 必有 $f(x_1) > f(x) > f(x_2)$, 则 $f(x)$ 在 $\mathbb{R}$ 上是严格单调减少函数.\par
% p048-b01; source PDF page 48
\par\Needspace{6\baselineskip}
\blockanchor{p048-b01}
\question{6}设 $f(x)$ 是 $(0, +\infty)$ 上的一致连续函数, $\{\lambda_n\}$ 是正数列, 若对任意固定的 $x \in (0, +\infty)$, 有\par
\nopagebreak[4]
% p048-b02; source PDF page 48
\blockanchor{p048-b02}
\displayformula{\lim _ {n \to \infty} f (x + \lambda_ {n}) = 0.}
% p048-b03; source PDF page 48
\blockanchor{p048-b03}
证明：函数列 $\{f(x + \lambda_n)\}$ 在闭区间 $[a,b](a > 0)$ 上一致收敛于零.\par
% p048-b04; source PDF page 48
\par\Needspace{6\baselineskip}
\blockanchor{p048-b04}
\question{7}设 $f(x)$ 是 $\mathbb{R}$ 上有界的连续函数, 若存在正常数 $a$, 使得\par
\nopagebreak[4]
% p048-b05; source PDF page 48
\blockanchor{p048-b05}
\displayformula{f (x) + a \int_ {x - 1} ^ {x} f (t) \mathrm{d} t}
% p048-b06; source PDF page 48
\blockanchor{p048-b06}
为 $\mathbb{R}$ 上的常值函数, 证明: $f(x)$ 是 $\mathbb{R}$ 上的常值函数.\par
